\subsection{Hamiltoniane}
$$
[\mathcal{H}_i(\hat{x}^{(i)}, \hat{p}^{(i)}), \mathcal{H}_j(\hat{x}^{(j)}, \hat{p}^{(j)})] = 0, \quad \text{se } i \neq j.
$$
La $ \mathcal{H} $ \`e somma di hamiltoniane che commutano fra loro, quindi gli autostati della $ \mathcal{H} $ sono tutti gli autostati di $ \mathcal{H}_i $. Allora
$$
\mathcal{H} \ket{\psi_E} = E \ket{\psi_E} \quad \Rightarrow \quad E_{k_1, \dots, k_d} = E_{k_1} + \cdots + E_{k_d}.
$$
Inoltre,
$$
\mathcal{H}_i \ket{\psi_k^{(i)}} = E_k \ket{\psi_k^{(i)}}, \qquad \braket{x^{(i)}|\psi_k^{(i)}} = \psi_k^{(i)}(x^{(i)})
$$
e ne segue che 
$$
\psi(\vec{x}) = \psi_{k_1}(x^{(1)}) \cdots \psi_{k_d}(x^{(d)}).
$$
Se il problema \`e separabile, ho questa utile condizione.

Esempio: 
$$
V(\vec{x}) = \sum_{i = 1}^3 V_i(x_i), \qquad V_i(x_i) = \begin{cases}
0 & |x_i| < a_i \\
+\infty & |x_i| > a_i
\end{cases}.
$$
Il problema si riduce a tre problemi unidimensionali.
$$
\left( \frac{{\hat{p}^{(i)}}^2}{2m} + V_i(x^i) \right) \ket{\psi_{n_i}} = E_{n_i} \ket{\psi_{n_i}},
$$
$$
E_{n_i} = \frac{\hbar^2 k_{n_i}^2}{2m}, \qquad k_{n_i} = n_i \frac{\pi}{2a_i},
$$
$$
\psi_{n_i}(x^{(i)}) = \begin{cases}
A_{n_i} \sin(k_{n_i} x_i) \; n_i \text{ pari} \\
B_{n_i} \cos(k_{n_i} x_i) \; n_i \text{ dispari}
\end{cases}.
$$

Nel caso di buca cubica, $ a_1 = a_2 = a_3 = a $ e lo spettro
$$
E_{n_1, n_2, n_3} = \frac{\hbar^2 \pi^2}{8ma^2} (n_1^2 + n_2^2 + n_3^2),
$$
$$
\text{stato fonamentale} \, E_{111} = \frac{3}{8} \frac{\hbar^2 \pi^2}{ma^2},
$$
$$
1^{\text{o}} \text{ stato eccitato (tre volte degenere) } E_{211} = E_{121} = E_{112} = \frac{3}{4} \frac{\hbar^2 \pi^2}{ma^2}.
$$